\section{REVISÃO BIBLIOGRÁFICA}
Este TCC investiga um modelo híbrido para
diagnóstico preditivo de falhas incipientes em transformadores de
potência a partir de Análise de Gases Dissolvidos (DGA).
O pipeline proposto possui dois estágios:
(i) um módulo de previsão temporal baseado em LSTM para estimar as
concentrações futuras dos gases; e
(ii) um módulo de classificação baseado em XGBoost, com seleção de atributos e otimização de hiperparâmetros via Algoritmo Genético Melhorado (IGA).
A fundamentação e nomenclatura de DGA e falhas seguem recomendações normativas e guias técnicos
\cite{iec60599_2022,ieee_c57_104_2019,cigre_tb771_2019}.
Aspectos de monitoramento online e confiabilidade de medições são
contextualizados por recomendações sobre monitores e verificação de acurácia \cite{cigre_tb783_2019}.
A motivação para o uso de ML e híbridos apoia-se em resultados
recentes com IGA+XGBoost \cite{wu2021_iga_xgboost} e em abordagens
dinâmicas com LSTM para previsão de falhas \cite{wang2023_lstm_dga}.
\textbf{Sugestão de figura:} inserir na Figura~\ref{fig:pipeline} um
diagrama do fluxo completo (DGA $\rightarrow$ pré-processamento $
\rightarrow$ LSTM $\rightarrow$ atributos derivados $\rightarrow$ IGA $
\rightarrow$ XGBoost $\rightarrow$ avaliação).

\section{Bases de dados e definição do problema}
\subsection{Origem e descrição dos dados}
Os dados consistem em séries temporais de medições de DGA
(concentrações de gases dissolvidos no óleo isolante), associadas a
rótulos de condição operacional e/ou classes de falha.
Nesta etapa de escrita, os seguintes itens estão \textbf{não
especificados} (devem ser preenchidos quando o dataset final for definido):
número de transformadores, intervalo temporal total coberto, periodicidade
de amostragem (regular vs.\ irregular), lista final de gases disponíveis,
unidade de medida (ppm ou $\mu$L/L), origem dos rótulos (inspeção,
equipe de manutenção, regra normatizada etc.), e proporção de
classes.

\subsection{Classes e rotulagem}
As classes podem seguir a taxonomia tradicional utilizada em DGA (por
exemplo, classes associadas a descargas e falhas térmicas) conforme guias
normativos \cite{iec60599_2022,ieee_c57_104_2019} e práticas discutidas em
revisões \cite{nanfak2024_traditional_dga_review,duval2002_review}.
Caso o dataset utilize classes diferentes (mais agregadas ou mais finas), o
mapeamento será documentado (\textbf{não especificado}).

\subsection{Estrutura do problema}
O problema é formulado como:
\begin{itemize}
\item \textbf{Previsão (regressão temporal):} prever concentrações futuras de gases (multi-saída) a partir de sequências históricas
\cite{wang2023_lstm_dga}.
\item \textbf{Classificação (multi-classe ou binária):} inferir
classe de falha/condição utilizando medições atuais, atributos
derivados e previsões \cite{wu2021_iga_xgboost}.
\end{itemize}

\section{Pré-processamento de dados}
\subsection{Limpeza, consistência e outliers}
Serão aplicados procedimentos para:
(i) eliminar valores impossíveis (negativos); (ii) ordenar por tempo; (iii) uniformizar unidade e escala (\textbf{não especificado}); e
(iv) tratar outliers por critérios robustos (IQR, z-score robusto) e/ou
regras de engenharia (\textbf{não especificado}).
A decisão final será reportada e justificada.

\subsection{Dados faltantes e imputação}
Dados faltantes são comuns em DGA (coleta laboratorial, falhas de
aquisição, integração de fontes).
Serão avaliadas estratégias de imputação:
interpolação temporal para lacunas curtas; \textit{forward fill} com
limite de tempo; e imputação multivariada (ex.: KNN/MICE).
Como linha de base orientada por literatura, pode-se considerar regressão
por Random Forest para preenchimento, como utilizado em diagnóstico com
XGBoost em cenários de dados incompletos
\cite{wang2023_tpe_xgboost_incomplete}.
A escolha final e hiperparâmetros permanecem \textbf{não especificados}.

\subsection{Normalização e transformações}
Embora XGBoost seja pouco sensível à escala, transformações podem
estabilizar variância e reduzir assimetria em gases.
Serão comparados:
$\log(1+x)$ por gás (quando aplicável) e normalização Min--Max ou z-score.
A configuração final será reportada (\textbf{não especificado}).

\subsection{Janelamento temporal (windowing)}
Para treinamento do LSTM, as séries temporais serão convertidas em
amostras supervisionadas via janela deslizante:
\begin{itemize}
\item comprimento da janela $L \in [4,48]$ passos (\textbf{não
especificado});
\item horizonte de previsão $H \in [1,12]$ passos (\textbf{não
especificado});
\item entrada $\mathbf{X}\in\mathbb{R}^{L\times m}$, com $m$
variáveis (gases e derivados);
\item saída $\mathbf{y}$ multi-saída (gases em $t+H$).
\end{itemize}
Para evitar vazamento temporal, splits e validação seguirão
cronologia (Seção~\ref{sec:protocolo}).
A literatura recente explicita que estratégias inadequadas podem
introduzir \textit{data leakage} em pipelines de previsão e decomposição, sugerindo amostragem passo-a-passo para maior realismo
\cite{chen2024_sds_vmd_lstm_dga}.

\section{Módulo de previsão: LSTM}
\subsection{Arquitetura candidata}
A rede LSTM será treinada para prever concentrações futuras de gases
a partir do histórico.
Arquiteturas candidatas (faixas):
\begin{itemize}
\item 1 a 2 camadas LSTM empilhadas;
\item unidades por camada $\in\{16,32,64,128\}$;
\item \textit{dropout} $\in[0.0,0.5]$ (incluindo \textit{recurrent
dropout} quando disponível);
\item camada densa final com saída multi-gás.
\end{itemize}
O desenho segue evidência de previsão dinâmica com LSTM aplicada à
DGA \cite{wang2023_lstm_dga}.

\subsection{Treinamento e função de perda}
Configurações iniciais e ranges:
\begin{itemize}
\item perdas: MSE, MAE e Huber;
\item otimizador: Adam;
\item taxa de aprendizado $\eta \in [10^{-5},10^{-2}]$;
\item \textit{batch size} $\in[16,256]$;
\item \'{e}pocas: até 500 com \textit{early stopping} (paciência
10--30).
\end{itemize}
Rotinas como redução de LR em platô e regularização L2 podem
ser testadas (\textbf{não especificado}).
A calibração final será documentada.

\subsection{Métricas do módulo de previsão}
Serão reportadas MAE e RMSE por gás e agregadas.
MAPE será usado apenas quando a escala permitir (tratando zeros).

\section{Engenharia de atributos para diagnóstico}
\subsection{Atributos derivados (razões e Duval)}
A partir de gases observados e previstos, serão extraídos:
\begin{itemize}
\item razões clássicas (incluindo Rogers/IEC quando aplicável) e
variáveis agregadas para caracterizar padrões de falha
\cite{iec60599_2022,ieee_c57_104_2019,nanfak2024_traditional_dga_review};
\item \textit{features} do
Triângulo de Duval: percentuais normalizados de CH$_4$, C$_2$H$_4$,
C$_2$H$_2$ e/ou codificação de zona como atributo categórico
\cite{duval2002_review,cigre_tb771_2019}.
\end{itemize}

\subsection{Uso de previsões como entradas do classificador}
O classificador receberá como entrada:
gases em $t$, previsões $\hat{\mathbf{g}}(t+H)$ e atributos derivados
calculados sobre $\hat{\mathbf{g}}$.
Esse desenho operacionaliza diagnóstico preditivo (antecipado) e segue a
motivação de previsão dinâmica proposta na literatura
\cite{wang2023_lstm_dga}.

10\section{IGA: seleção de atributos e ajuste de hiperparâmetros}
\subsection{Codificação do cromossomo}
O cromossomo do IGA será composto por:
(i) uma parte binária (máscara) para selecionar atributos; e
(ii) genes inteiros/reais para hiperparâmetros do XGBoost.
Esse tipo de otimização conjunta foi utilizado com sucesso em
diagnóstico por DGA com XGBoost \cite{wu2021_iga_xgboost}.

\subsection{Função de aptidão (fitness) e restrições}
Como métrica principal recomenda-se Macro-F1 (robusta a desbalanceamento),
com penalização proporcional ao número de atributos selecionados:
$$
F = \alpha\,(1-\text{MacroF1}) + (1-\alpha)\,\frac{k}{d},
$$
onde $k$ é o número de atributos selecionados, $d$ é o total de
atributos e $\alpha \in [0.7,0.95]$ (\textbf{não especificado}).
Restrições sugeridas:
\begin{itemize}
\item $k \ge k_{\min}$ para evitar modelos degenerados (\textbf{não
especificado});
\item limites de hiperparâmetros do XGBoost (Seção seguinte).
\end{itemize}

\subsection{Operadores e parâmetros do IGA}
Ranges iniciais:
população 50--200; gerações 50--200; seleção por torneio +
elitismo (1--5\%);
cruzamento 0.6--0.9; mutação 0.01--0.2.
A configuração final será reportada (\textbf{não especificado}) e
pode ser calibrada com base em desempenho e tempo computacional
\cite{wu2021_iga_xgboost}.

\section{Classificador XGBoost}
\subsection{Justificativa e objetivo}
O XGBoost é um sistema de \textit{gradient boosting} em árvores
amplamente consolidado para dados tabulares, conhecido por escalabilidade e
desempenho \cite{chen2016_xgboost}.
No presente TCC, ele será empregado como classificador de classes de falha
(multi-classe) ou condição (binária), conforme o dataset final
(\textbf{não especificado}).

\subsection{Configuração}
Configuração sugerida:
\begin{itemize}
\item multi-classe: \texttt{objective=multi:softprob} e
\texttt{eval\_metric=mlogloss};
\item binária: \texttt{objective=binary:logistic} e
\texttt{eval\_metric=logloss}.
\end{itemize}
11Quando possível, utilizar \textit{early stopping} com validação
temporal.

\subsection{Ranges de hiperparâmetros (para busca via IGA)}
Faixas recomendadas:
\begin{itemize}
\item \texttt{n\_estimators} $\in[200,3000]$
\item \texttt{max\_depth} $\in[2,12]$
\item \texttt{learning\_rate} $\in[0.005,0.3]$
\item \texttt{subsample} $\in[0.5,1.0]$
\item \texttt{colsample\_bytree} $\in[0.5,1.0]$
\item \texttt{min\_child\_weight} $\in[1,20]$
\item \texttt{gamma} $\in[0,10]$
\item \texttt{reg\_alpha} $\in[0,10]$
\item \texttt{reg\_lambda} $\in[0.1,50]$
\end{itemize}

\section{Protocolo experimental e avaliação}\label{sec:protocolo}
\subsection{Split treino/validação/teste}
Por se tratar de série temporal, o split seguirá ordem cronológica:
treino (período antigo), validação (intermediário) e teste (final).
Opcionalmente, usar \textit{time-series cross-validation} do tipo
\textit{rolling-origin} (expanding window) quando houver dados suficientes
(\textbf{não especificado}).
A necessidade de evitar vazamento temporal é coerente com preocupações recentes em previsão de gases \cite{chen2024_sds_vmd_lstm_dga}.

\subsection{Métricas de classificação}
Serão reportadas:
acurácia; precisão, revocação e F1 (macro e ponderada);
AUC One-vs-Rest (quando aplicável);
e matriz de confusão.

\subsection{Baselines para comparação}
Baselines:
\begin{itemize}
\item métodos tradicionais: Duval e razões (incluindo Rogers/IEC
quando aplicável), conforme guias e revisões
\cite{iec60599_2022,ieee_c57_104_2019,cigre_tb771_2019,nanfak2024_traditional_dga_review,duval2002_review}
\item ML estático: SVM e Random Forest (\textbf{configurações:
não especificado});
\item XGBoost sem previsão e sem IGA (baseline de boosting).
\end{itemize}

\subsection{Plano de ablação}
Para quantificar a contribuição de cada componente:
\begin{itemize}
\item (A) XGBoost apenas com gases atuais;
\item (B) LSTM + XGBoost (sem IGA);
\item (C) IGA + XGBoost (sem previsão);
\item (D) LSTM + IGA + XGBoost (completo);
12\item (E) remoção seletiva de \textit{features} (razões vs.\
Duval vs.\ taxas).
\end{itemize}

\section{Implementação, recursos computacionais e reprodutibilidade}
\subsection{Ferramentas}
Sugere-se Python com \texttt{numpy}, \texttt{pandas}, \texttt{scikit-learn},
\texttt{xgboost} e um framework DL (\texttt{TensorFlow/Keras} ou
\texttt{PyTorch}).
Versões exatas: \textbf{não especificado}.

\subsection{Hardware e controle de aleatoriedade}
Reportar CPU/GPU, RAM e tempo de treinamento (\textbf{não especificado}).
Fixar sementes (Python/NumPy/DL), versionar código (Git) e registrar
experimentos (ex.: MLflow) (\textbf{não especificado}).

\section{Limitações e extensões (tendências)}
\subsection{Monitoramento online/IoT}
A disponibilidade crescente de monitores online e a necessidade de requisitos
de acurácia em contratos e aquisições são discutidas por
relatórios técnicos \cite{cigre_tb783_2019}.
Como exemplo de tendência de integração com IoT e métodos
inteligentes, trabalhos propõem classificação de falhas com
sensoriamento online e modelos neuro-fuzzy \cite{equbal2022_iot_anfis_enose}.

\subsection{Mudança de distribuição e adaptação de domínio}
Modelos treinados em uma base podem degradar sob condições operacionais diferentes (\textit{domain shift}).
Como extensão, recomenda-se avaliar adaptação de domínio com
ponderação de atributos e alinhamento de distribuição
\cite{mahmoodiyan2025_domain_adaptation_thesis}.

\subsection{Aprendizado federado/transferência federada}
Em cenários multi-organização com restrições de privacidade e
\textit{data islands}, abordagens de transferência federada são
relevantes para lidar com discrepância de distribuição e custo de
rotulagem \cite{qian2025_ftl_review}.
(C) Tabela comparativa e lacunas
\subsection{Energia}
revisão bibliográfica sobre energia...